% appP_graph_a.tex -- generated by writing/tools/blueprint_gen.py 1.0 (level 3) on 2026-10-14 from sources a9a6d40381; do not edit
\providecommand{\bpNum}[1]{\hyperref[#1]{\ref*{#1}}}%
\providecommand{\bpNumA}[1]{\hyperref[#1]{\textcolor{orange!65!black}{\ref*{#1}}}}%
\providecommand{\bpPage}[1]{\pageref*{#1}}%
\providecommand{\bpLab}[1]{{\ttfamily #1}}%
\providecommand{\bpBrk}{\allowbreak}%
\providecommand{\bpT}[2]{\facttagbox[#1]{\scriptsize #2}}%
\providecommand{\bpU}[1]{{\color{gray!70!black}#1}}%
\section{依赖图}\label{sec:appP-graphs}
本节的图由 \texttt{\string\uses} 数据自动生成：\cref{fig:appP-dependency}为全书骨架，其余各图依次给出各部分的主要定理（每图至多 40 个节点，超出时按折叠规则省略使用者最少的次要前提）。
\begin{figure}[p]
\centering
\rotatebox{90}{\begin{tikzpicture}[
  bp node/.style={draw,line width=0.45pt,rounded corners=1pt,font=\tiny\ttfamily,align=center,
                  inner xsep=1.5pt,inner ysep=1.1pt,minimum height=4mm},
  bp R/.style={bp node,fill=OliveGreen!15,draw=OliveGreen!60!black},
  bp A/.style={bp node,fill=RoyalBlue!15,draw=RoyalBlue!60!black},
  bp P/.style={bp node,fill=BurntOrange!15,draw=BurntOrange!70!black},
  bp C/.style={bp node,fill=Plum!15,draw=Plum!70!black},
  bp V/.style={bp node,fill=BrickRed!15,draw=BrickRed!70!black,double,double distance=0.6pt},
  bp cond/.style={dash pattern=on 1.6pt off 1pt},
  bp math/.style={bp node,fill=white,draw=black!60},
  bp def/.style={bp node,fill=white,draw=black!45,line width=0.3pt,rounded corners=0pt},
  bp conj/.style={bp node,ellipse,inner xsep=0pt,inner ysep=0pt,fill=Plum!15,draw=Plum!70!black},
  bp asm/.style={bp node,fill=orange!8,draw=orange!65!black,dash pattern=on 2pt off 1.2pt},
  bp stub/.style={bp node,font=\tiny,fill=gray!12,draw=gray!55,text=black!70},
  bp fold/.style={bp node,font=\tiny,fill=white,draw=gray!55,dash pattern=on 1pt off 1pt,text=black!70},
  bp unit/.style={font=\scriptsize,text=black!70,inner sep=1pt},
  bp sep/.style={draw=gray!40,line width=0.4pt,dash pattern=on 1pt off 1.5pt},
  bp key/.style={font=\scriptsize,anchor=west,inner sep=0pt},
  bp use/.style={draw=gray!60,line width=0.5pt,-{Stealth[length=2.6pt,width=2.2pt]}},
  bp stubuse/.style={draw=gray!45,line width=0.45pt,dash pattern=on 1.2pt off 1pt,-{Stealth[length=2.6pt,width=2.2pt]}},
  bp asmuse/.style={draw=orange!65!black,line width=0.5pt,dash pattern=on 2pt off 1.2pt,-{Stealth[length=2.6pt,width=2.2pt]}},
  bp neck/.style={draw=red!60!black,line width=1.1pt,-{Stealth[length=3.6pt,width=3.2pt]}},
  bp long/.style={opacity=0.5}]
  \node[bp unit] at (0.738,3.574) {第~1--5~章};
  \node[bp unit] at (4.390,3.574) {第~6--8~章};
  \node[bp unit] at (7.785,3.574) {第~9--13~章};
  \node[bp unit] at (11.004,3.574) {第~14--18~章};
  \node[bp unit] at (13.234,3.574) {第~19--20~章};
  \node[bp unit] at (15.472,3.574) {第~21~章};
  \node[bp unit] at (18.382,3.574) {第~22--24~章};
  \draw[bp sep] (2.437,3.434) -- (2.437,-3.304);
  \draw[bp sep] (6.097,3.434) -- (6.097,-3.304);
  \draw[bp sep] (9.622,3.434) -- (9.622,-3.304);
  \draw[bp sep] (12.156,3.434) -- (12.156,-3.304);
  \draw[bp sep] (14.312,3.434) -- (14.312,-3.304);
  \draw[bp sep] (16.631,3.434) -- (16.631,-3.304);
  \node[bp R] (n1) at (0.000,1.521) {\detokenize{ir-}\\\detokenize{landau-}\\\detokenize{momentum}};
  \node[bp R] (n2) at (0.000,0.383) {\detokenize{ir-}\\\detokenize{power-}\\\detokenize{counting}};
  \node[bp A] (n3) at (1.475,0.862) {\detokenize{ir-}\\\detokenize{leading-}\\\detokenize{regions}};
  \node[bp A] (n4) at (0.000,-0.510) {\detokenize{ir-kln}};
  \node[bp A] (n5) at (1.475,2.036) {\detokenize{ir-safe-}\\\detokenize{finite}};
  \node[bp R] (n6) at (0.000,-1.404) {\detokenize{ir-}\\\detokenize{regions-}\\\detokenize{euclid}};
  \node[bp A] (n7) at (1.475,-1.297) {\detokenize{sudakov-}\\\detokenize{factorization}};
  \node[bp A] (n8) at (1.475,-2.188) {\detokenize{sudakov-}\\\detokenize{ir-poles}};
  \node[bp R] (n9) at (3.514,2.271) {\detokenize{ee-}\\\detokenize{euclid-}\\\detokenize{ir}};
  \node[bp R] (n10) at (5.265,2.795) {\detokenize{ee-R-}\\\detokenize{finite}};
  \node[bp conj] (n11) at (5.265,1.808) {\detokenize{ee-}\\\detokenize{duality}};
  \node[bp R] (n12) at (3.514,0.697) {\detokenize{pdf-}\\\detokenize{support-}\\\detokenize{pt}};
  \node[bp R] (n13) at (5.265,0.697) {\detokenize{pdf-}\\\detokenize{moments-}\\\detokenize{local}};
  \node[bp R] (n14) at (5.265,-0.317) {\detokenize{pdf-}\\\detokenize{dglap}};
  \node[bp R] (n15) at (5.265,-1.549) {\detokenize{pdf-dis-}\\\detokenize{xspace}};
  \node[bp conj] (n16) at (3.514,-0.712) {\detokenize{pdf-}\\\detokenize{positivity}};
  \node[bp A] (n17) at (6.944,1.205) {\detokenize{method-}\\\detokenize{subtraction}};
  \node[bp A] (n18) at (6.944,0.191) {\detokenize{method-}\\\detokenize{glauber-}\\\detokenize{deformation}};
  \node[bp A] (n19) at (8.626,0.500) {\detokenize{method-}\\\detokenize{template}};
  \node[bp A] (n20) at (8.626,2.933) {\detokenize{disall-}\\\detokenize{factorization}};
  \node[bp A] (n21) at (8.626,2.042) {\detokenize{disall-}\\\detokenize{heavy}};
  \node[bp A] (n22) at (8.626,-0.392) {\detokenize{dy-}\\\detokenize{factorization}};
  \node[bp A] (n23) at (8.626,-1.406) {\detokenize{dy-}\\\detokenize{colour-}\\\detokenize{singlet}};
  \node[bp conj] (n24) at (8.626,-2.677) {\detokenize{dy-}\\\detokenize{coloured-}\\\detokenize{final}};
  \node[bp R] (n25) at (10.411,1.105) {\detokenize{scet-}\\\detokenize{decoupling}};
  \node[bp A] (n26) at (11.598,0.687) {\detokenize{scet-}\\\detokenize{thrust}};
  \node[bp A] (n27) at (10.411,0.068) {\detokenize{jets-bms}};
  \node[bp A] (n28) at (13.234,0.807) {\detokenize{heavy-}\\\detokenize{bdpi}};
  \node[bp A] (n29) at (13.234,1.821) {\detokenize{heavy-}\\\detokenize{inclusive-}\\\detokenize{b}};
  \node[bp A] (n30) at (13.234,-0.208) {\detokenize{heavy-}\\\detokenize{nrqcd-}\\\detokenize{decay}};
  \node[bp conj] (n31) at (13.234,-1.478) {\detokenize{heavy-}\\\detokenize{nrqcd-}\\\detokenize{production}};
  \node[bp conj] (n32) at (15.472,2.211) {\detokenize{power-}\\\detokenize{renormalon-}\\\detokenize{ht}};
  \node[bp R] (n33) at (15.472,-0.182) {\detokenize{power-}\\\detokenize{ope-}\\\detokenize{convergence}};
  \node[bp conj] (n34) at (17.627,1.955) {\detokenize{lamet-}\\\detokenize{principle}};
  \node[bp A] (n35) at (17.627,-0.300) {\detokenize{lattice-}\\\detokenize{lamet}};
  \node[bp asm] (n36) at (17.627,0.967) {\detokenize{lamet-}\\\detokenize{staple}};
  \node[bp asm] (n37) at (19.137,1.536) {\detokenize{lamet-}\\\detokenize{order-}\\\detokenize{limits}};
  \node[bp A] (n38) at (19.137,0.399) {\detokenize{lamet-}\\\detokenize{quasi-}\\\detokenize{tmd}};
  \node[bp asm] (n39) at (17.627,-1.314) {\detokenize{yzliu-}\\\detokenize{wightman-}\\\detokenize{pt}};
  \node[bp A] (n40) at (19.137,-1.308) {\detokenize{yzliu-}\\\detokenize{time-}\\\detokenize{ordering}};
  \begin{scope}[on background layer]
  \draw[bp use] (n20) to[bend left=35] (n21);
  \draw[bp use] (n22) to[bend left=35] (n23);
  \draw[bp use] (n22) to[bend left=35] (n24);
  \draw[bp use] (n25) -- (n26);
  \draw[bp use] (n25) -- (n28);
  \draw[bp use] (n25) -- (n29);
  \draw[bp use] (n36) -- (n37);
  \draw[bp use] (n1) -- (n3);
  \draw[bp use] (n2) -- (n3);
  \draw[bp use] (n3) to[bend left=35] (n5);
  \draw[bp use] (n3) to[bend left=35] (n7);
  \draw[bp use,bp long] (n3) -- (n17);
  \draw[bp use,bp long] (n3) -- (n18);
  \draw[bp use] (n6) -- (n7);
  \draw[bp use] (n7) to[bend left=35] (n8);
  \draw[bp use,bp long] (n9) -- (n32);
  \draw[bp use] (n9) -- (n10);
  \draw[bp use] (n9) -- (n11);
  \draw[bp use] (n12) -- (n13);
  \draw[bp use] (n13) to[bend left=35] (n14);
  \draw[bp use,bp long] (n14) -- (n35);
  \draw[bp use] (n14) to[bend left=35] (n15);
  \draw[bp use] (n14) -- (n19);
  \draw[bp use] (n17) -- (n19);
  \draw[bp use,bp long] (n18) -- (n26);
  \draw[bp use] (n18) -- (n19);
  \draw[bp use] (n19) to[bend left=35] (n20);
  \draw[bp use] (n19) to[bend left=35] (n22);
  \draw[bp neck] (n36) -- (n38);
  \draw[bp neck] (n37) to[bend left=35] (n38);
  \draw[bp neck] (n39) -- (n40);
  \draw[bp neck] (n11) -- (n29);
  \end{scope}
  \begin{scope}[shift={(-0.524,-3.804)}]
    \node[bp R,anchor=west,minimum width=4.5mm] (k1) at (0,0) {R};
    \node[bp R,bp cond,anchor=west,minimum width=4.5mm,right=0.8mm of k1] (k2) {R};
    \node[bp A,anchor=west,minimum width=4.5mm,right=0.8mm of k2] (k3) {A};
    \node[bp P,anchor=west,minimum width=4.5mm,right=0.8mm of k3] (k4) {P};
    \node[bp C,anchor=west,minimum width=4.5mm,right=0.8mm of k4] (k5) {C};
    \node[bp V,anchor=west,minimum width=4.5mm,right=0.8mm of k5] (k6) {V};
    \node[bp key,right=1.2mm of k6] (t1) {严格性标签（虚线框为 $R_{\rm cond}$）};
    \node[bp conj,anchor=west,minimum width=7mm,minimum height=3.4mm,right=2.5mm of t1] (k7) {};
    \node[bp key,right=1mm of k7] (t2) {猜想};
    \node[bp asm,anchor=west,minimum width=4.5mm,right=2.5mm of t2] (k8) {};
    \node[bp key,right=1mm of k8] (t3) {编号假设};
    \node[bp math,anchor=west,minimum width=4.5mm,right=2.5mm of t3] (k9) {};
    \node[bp key,right=1mm of k9] (t4) {纯数学};
    \node[bp stub,anchor=west,minimum width=4.5mm,right=2.5mm of t4] (k10) {};
    \node[bp key,right=1mm of k10] (t5) {其他单元};
    \node[bp fold,anchor=west,minimum width=4.5mm,right=2.5mm of t5] (k11) {$+k$};
    \node[bp key,right=1mm of k11] {折叠};
    \draw[bp use] (0,-0.42) -- ++(6mm,0) coordinate (e1);
    \node[bp key] (u1) at ($(e1)+(1.2mm,0)$) {依赖（前提指向陈述）};
    \draw[bp asmuse] ($(u1.east)+(2.5mm,0)$) -- ++(6mm,0) coordinate (e2);
    \node[bp key] (u2) at ($(e2)+(1.2mm,0)$) {编号假设 $\to R_{\rm cond}$};
    \draw[bp neck] ($(u2.east)+(2.5mm,0)$) -- ++(6mm,0) coordinate (e3);
    \node[bp key] (u3) at ($(e3)+(1.2mm,0)$) {瓶颈依赖（主表）};
    \draw[bp use,bp long] ($(u3.east)+(2.5mm,0)$) -- ++(6mm,0) coordinate (e4);
    \node[bp key] (u4) at ($(e4)+(1.2mm,0)$) {跨两列以上};
    \draw[bp stubuse] ($(u4.east)+(2.5mm,0)$) -- ++(6mm,0) coordinate (e5);
    \node[bp key] at ($(e5)+(1.2mm,0)$) {来自其他单元};
  \end{scope}
\end{tikzpicture}}
\caption{本书主要陈述的依赖图（骨架；由各环境中的 \texttt{\string\uses} 自动生成）：总表中的主要陈述、它们的瓶颈假设与附录~M--O 的 $R_{\rm cond}$ 定理（40 个节点，另有 82 个陈述按折叠规则未画出；按书序自左向右分层，每层为一部分，层内按依赖深度分列；其余各图每层为一个单元）。框的颜色为严格性标签，虚线框为 $R_{\rm cond}$ 条件定理，橙色虚线框为编号假设，椭圆为猜想，双线框为 [V]，白色框为纯数学。边由前提指向陈述，是可达关系在所画节点上的传递约化（经由未画出陈述的依赖也画成一条边）；橙色虚线边为编号假设到 $R_{\rm cond}$ 陈述的依赖，粗红线为决定全阶论证地位的瓶颈依赖（取自主表“瓶颈”列），半透明线跨越两列以上。}\label{fig:appP-dependency}
\end{figure}
